% !PRINTED 88
\weberpage{88}

To obtain a system of necessary, sufficient, and mutually independent conditions, we state the question somewhat more precisely and ask for the conditions:

\gesperrt{that from a system of \(m\) linear homogeneous equations with \(n\) unknowns, \(\nu\) of the unknowns can be completely determined by \(n-\nu\) which remain arbitrary.}

This question too is in fact already answered in \S.\,24. Among the \(\nu\)-rowed determinants one must be different from zero, while all the \((\nu+1)\)-rowed determinants vanish. But it is already sufficient if it is known of a smaller number of the \((\nu+1)\)-rowed determinants that they vanish.

Assume that the unknowns \(x_{\nu+1},x_{\nu+2},\ldots,x_n\) are to remain arbitrary, that \(x_1,x_2,\ldots,x_\nu\) are to be determined by them, and take the determinant
\begin{equation}
        A=\begin{vmatrix}
        a_1^{(1)}&a_2^{(1)}&\cdots&a_\nu^{(1)}\\
        a_1^{(2)}&a_2^{(2)}&\cdots&a_\nu^{(2)}\\
        \vdots&\vdots&&\vdots\\
        a_1^{(\nu)}&a_2^{(\nu)}&\cdots&a_\nu^{(\nu)}
        \end{vmatrix}
\tag{1}
\end{equation}
to be different from zero. We calculate the unknowns \(x_1,x_2,\ldots,x_\nu\) according to \S.\,24, (9), and form the sums (11), whose vanishing says that the given system of equations is actually satisfied. The conditions for this are therefore:
\begin{equation}
\begin{vmatrix}
        a_1^{(1)}&a_2^{(1)}&\cdots&a_\nu^{(1)}&a_h^{(1)}\\
        a_1^{(2)}&a_2^{(2)}&\cdots&a_\nu^{(2)}&a_h^{(2)}\\
        \vdots&\vdots&&\vdots&\vdots\\
        a_1^{(\nu)}&a_2^{(\nu)}&\cdots&a_\nu^{(\nu)}&a_h^{(\nu)}\\
        a_1^{(k)}&a_2^{(k)}&\cdots&a_\nu^{(k)}&a_h^{(k)}
\end{vmatrix}=0,
\tag{2}
\end{equation}
or, written otherwise,
\begin{equation}
        A a_h^{(k)}-\sum_{i=1}^{\nu}\sum_{\mu=1}^{\nu}
        a_h^{(i)}a_\mu^{(k)}A_\mu^{(i)}=0,
\tag{3}
\end{equation}
and these conditions are also sufficient. Equation (2) or (3), however, is identically satisfied when \(h=1,2,\ldots,\nu\), or \(k=1,2,\ldots,\nu\), and therefore gives no condition on the coefficients for these values. Such conditions arise only for
%CONT>>
